\begin{textsample}{POS Dim 9 – global\_south\_2023\_11 – Score 3.00 – global\_south\_2023\_11/103122405.txt}  \label{ex:f9_pos_127}
The global cyber divide between Gaza and Israel Nov 1, 2023 The conflict between Gaza and Israel has intensified to new heights over the past few weeks, gaining global attention unlike any other.
This escalation of tensions has brought about a new battle line, one that only in recent years has become a pillar to modern-day warfare.
By Patrick McAteer, cyber threat intelligence analyst at SecurityHQ While the bulk of the fighting has taking place on the ground and been kept closed off to the region, the cyber war has begun to stir and expand worldwide.
For years, nations have been putting one another ' s cyber defence systems to the test.
The difference here is the mass international involvement from groups located far outside the region joining the digital fight.
Silent sabotage Similar to the further escalation of the Russia-Ukraine war that began almost two years ago, threat groups, individuals and other APT ' s have used this chance to exploit various realms across the cyber field.
Each with their own motive for joining the fight and idea of @ @ @ @ @ @ @ @ @ @ Some join for their own personal amusement, others for monetary gains, but the majority have joined for the political reasons embedded in this deep-rooted conflict.
Just days after Hamas entered Israel on 7 October, we observed a significant increase in the number of threat actors that began to align with one side or the other which, in turn, has also caused a significant increase in the number of cyber-attacks occurring in the region and around the world.
As of right now, there have only been a few serious cases, with the majority causing minimal damage.
There are however a few possible developments which could lead to heavier consequences than what we have seen so far, with Russian, Iranian, \textbf{Indian}, and other nation-backed groups recently entering the fray.
Disruption dilemma Since the outbreak of the current conflict, DDoS ( Distributed Denial of Service ) has overwhelmingly dominated the landscape of cyber-attack methods being used.
While most incidents have resulted in limited damage, we observed instances where overwhelming traffic managed to @ @ @ @ @ @ @ @ @ @, telecommunication, and critical infrastructure websites.
These, however, are not permanent solutions, as they last only several minutes or hours.
DDoS has not been the exclusive attack method throughout this conflict ; with cases of website defacement and general vulnerability exploits also taking place.
During the initial attack, a mobile app that alerts Israeli civilians of incoming rockets known as"Red Alert : Israel"had a vulnerability exploited.
This allowed the hackers to not only intercept alert requests, but also to dispatch fake ones including one that said,"nuclear bomb is coming".
With thousands more downloading the various warning \textbf{applications} after the drastic influx of rockets, it was also revealed that malicious versions posing as the official"RedAlert"\textbf{application} were being installed all over Israel.
After a few days they had obtained access to sensitive information from both iOS and Android stores belonging to over 100 000 users.
The sophisticated spyware within took code from the authentic RedAlert \textbf{application}, but requested further access to @ @ @ @ @ @ @ @ @ @, SMS Messages, list of Installed Software, Logged in Email, other App Accounts and more.
Attacks from sympathisers on both sides of the conflict have caused varied levels of disruption over the past few weeks.
For now, the damage has not escalated to much more than websites in both the public and private sector being taken down, defaced or various other exploits being taken advantage of.
With the number of groups continuing to expand, there seems to be a sudden surge in malicious activity which places an increased number of organisations at risk.
Development of threat groups Although the conflict is contained to a small section of the Levant region, the overwhelming mass mobilisation across the world has helped draw a much larger picture of the political landscape rapidly unfolding.
Cyber groups worldwide have pledged their services to varying sides, with the majority siding alongside the pro-Palestine or anti-Israel banner.
Many seem to be following the path that their nation takes, with pro-Russian groups like"Killnet"taking up digital @ @ @ @ @ @ @ @ @ @ nations like"Kerala Cyber Xtractors"in India hitting Palestine.
From state-backed threat actors to individual hacktivists, the wide variety of groups is increasing day by day and does not seem to be slowing down.
The battleground does not stop at the borders of just these two nations though, with cyber-attacks from either side targeting critical infrastructure, military targets, governments, and private organisations of countries that are aligning with the opposing side.
Nations with pro-Palestinian hacktivists have set their sights on nations such as US, UK, Germany, India, France, Italy, Canada and more over recent weeks.
Likewise, various pro-Israeli hacktivists have targeted, Iran, Iraq, Qatar, \textbf{Saudi} Arabia and Lebanon.
While the list of threat groups is ever changing, here are the current known groups that are actively participating, the side of the conflict they chose to align with and their preferred attack vector.

% matched lemmas: application, indian, saudi
\end{textsample}
